%======================================================================
% CAPÍTULO 05: NÓMINA DE PENSIONADOS
%======================================================================
% Cálculos, conceptos, retroactivos y calculadora
% Fundamento Legal: Art. 7, 34-47 de la Ley SSSIFANB
%======================================================================

\chapter{Módulo de Nómina de Pensionados}
\label{ch:nomina-pensionados}

\section{Descripción General}

Este módulo suple el Capítulo III (Art. 34 al 47) protegiendo el dictamen de que \textbf{todas las pensiones militares son inembargables (Art. 7)}.

El Módulo de Nómina de Pensionados gestiona el cálculo, procesamiento y disbursamiento de las pensiones de los beneficiarios del sistema.

\section{Submódulos Disponibles}

\begin{itemize}[leftmargin=*]
    \item \textbf{Cálculos}: Procesamiento de nómina mensual
    \item \textbf{Conceptos}: Tabulador de conceptos salariales
    \item \textbf{Patria}: Integración con Sistema Patria
    \item \textbf{Retroactivos}: Homologación de pensiones
    \item \textbf{Calculadora}: Simulador de pensiones
\end{itemize}

\section{Cálculos de Nómina - Reserva Activa (Art. 34 y 38)}

El sistema aplica la Matriz Ponderada de Pensiones de Reserva Activa:

\keyconcept{Matriz Ponderada de Pensiones de Reserva Activa (Mandato Art. 38)}\\
El controlador maestro restringe y asigna los porcentajes garantizando la retribución exacta por los años de servicio, validando topes máximos y mínimos legales.

\begin{lstlisting}[language=Java, caption=Algoritmo de Matriz Ponderada]
function generarPensionActiva(anosServicio, ultimoSueldoIntegral) {
    if (anosServicio < 15) {
        return { error: 'REJECT_INSUFFICIENT_TIME' }; // No califica
    }
    
    let porcentaje = 0;
    if (anosServicio === 15) porcentaje = 0.50;
    else if (anosServicio >= 30) porcentaje = 1.00;
    else {
        // Ponderacion fraccionaria intermedia entre 15..29
        porcentaje = calcularFactorIntermedio(anosServicio); 
    }
    
    let pensionCalculada = ultimoSueldoIntegral * porcentaje;
    
    // Mecanismo de Piso Economico (Art. 35 y 36)
    if (pensionCalculada < SALARIO_MINIMO_NACIONAL) {
        pensionCalculada = SALARIO_MINIMO_NACIONAL;
    }
    
    return pensionCalculada;
}
\end{lstlisting}

\subsection*{Proceso de Cálculo}

El sistema realiza mensualmente el cálculo de nómina considerando:

\begin{enumerate}[leftmargin=*]
    \item Identificación del pensionado activo
    \item Determinación del tipo de pensión
    \item Aplicación del tabulador de conceptos
    \item Cálculo de deducciones aplicables
    \item Determinación del neto a pagar
    \item Generación de orden de pago
\end{enumerate}

\begin{tcolorbox}[
    enhanced,
    colback=white,
    colframe=sandraDark,
    width=0.85\textwidth,
    halign=center
]
\begin{tabular}{lr}
\toprule
\textbf{Concepto} & \textbf{Monto} \\
\midrule
Pensión Base & \$\$5.000.000,00 \\
Complemento por Hijos & +\$\$500.000,00 \\
Asignación Especial & +\$\$200.000,00 \\
\midrule
\textbf{Bruto} & \textbf{\$\$5.700.000,00} \\
\midrule
Deducción ISLR & -\$\$285.000,00 \\
Deducción SSO (6.5\%) & -\$\$370.500,00 \\
\midrule
\textbf{Neto a Pagar} & \textbf{\$\$5.044.500,00} \\
\bottomrule
\end{tabular}
\end{tcolorbox}
\caption{Ejemplo de cálculo de nómina}
\label{fig:calculo-nomina}
\end{figure}

\section{Tabulador de Conceptos}

\subsection*{Conceptos de Ingreso}

\begin{longtable}{p{4cm}p{8cm}}
\toprule
\textbf{Concepto} & \textbf{Descripción} \\
\midrule
\endhead
\bottomrule
\caption{Conceptos de ingreso pensionado}
\label{tab:conceptos-ingreso}
\endlastfoot

Pensión Base & Monto según grado y años de servicio (Art. 38) \\
Complemento por Antigüedad & Incremento por años de servicio \\
Asignación por Hijos & Complemento por dependientes menores \\
Asignación Especial & Bono discrecional autorizado \\
Retroactivo & Diferencia pendiente de homologación (Art. 39) \\
Prima por Responsabilidad & Complemento por cargo especial \\
\end{longtable}

\subsection*{Conceptos de Deducción}

\begin{longtable}{p{4cm}p{8cm}}
\toprule
\textbf{Concepto} & \textbf{Descripción} \\
\midrule
\endhead
\bottomrule
\caption{Conceptos de deducción pensionado}
\label{tab:conceptos-deduccion}
\endlastfoot

ISLR & Impuesto sobre la renta (escala progresiva) \\
SSO (Salud) & Aporte de salud 6.5\% del monto pensión \\
LSRQ & Ley de Política de Vivienda \\
Embargo Judicial & Deducción por medida cautelar \\
Manutención & Pensión alimenticia ordenada \\
\end{longtable}

\section{Vector de Sobrevivientes e Invalidez (Art. 43 y 46)}

El núcleo que previene el desamparo tras el fallecimiento del servidor.

\subsection*{Vector de Distribución de Sobrevivientes (Art. 43)}

\begin{lstlisting}[language=Java, caption=Matriz de Sobrevivientes]
// El algoritmo divide legalmente el capital para evitar desamparo familiar
const liquidarSobrevivientes = (montoTotalFiduciario, sobrevivientesValid) => {
    let asignaciones = {
        conyuge: montoTotalFiduciario * 0.60,
        hijos: montoTotalFiduciario * 0.20,      // Total dividido entre num de hijos
        padres: montoTotalFiduciario * 0.20      // 10% Madre, 10% Padre
    };

    // Ruteo de Fallas Familiares (Exclusión de Supervivencia)
    if (!sobrevivientesValid.includes('padres')) {
        // Reubicacion matematica algoritmica sobre el remanente
        reubicarPorcentajeFaltante(asignaciones, 'padres', 0.20);
    }
    
    return asignaciones;
};
\end{lstlisting}

\subsection*{Generadores por Invalidez (Art. 46)}

\begin{tcolorbox}[
    colback=sandraSoft,
    colframe=sandraDark,
    width=\textwidth,
    title=\textbf{Invalidez en Acto de Servicio}
]
\small
En vida, el sistema genera nómina a militares en Acto de Servicio garantizando sus derechos prestacionales inmediatos:
\begin{itemize}
    \item \textbf{Invalidez Absoluta Permanente}: Se asigna el \textbf{100\%} de la tabla equivalente a la posición del retiro de su escalafón.
    \item \textbf{Invalidez Parcial Permanente}: Asume el \textbf{75\%} de los montos ante limitación que impida al servidor realizar labores regulares.
\end{itemize}
\end{tcolorbox}

\section{Integración Sistema Patria}

El sistema se integra con el Sistema Patria para la validación y transferencia de pagos:

\begin{tcolorbox}[
    enhanced,
    colback=sandraSoft,
    colframe=sandraDark,
    width=0.9\textwidth,
    halign=center
]
\small
La integración con el Sistema Patria permite verificar la condición de vida del pensionado, validar el método de pago y realizar transferencias directas a la billetera digital del beneficiario.
\end{tcolorbox}

\section{Retroactivos (Art. 39 - Homologación)}

Los retroactivos surgen cuando se producen homologaciones salariales:

\begin{tcolorbox}[
    enhanced,
    colback=white,
    colframe=sandraAccent,
    width=0.9\textwidth,
    halign=center
]
\small
\textbf{Cron Operativo Transversal}: Las actualizaciones ejecutadas por un analista en ``Configuración / Parámetros'' (Cambiar el tabulador Base Militar Activa por un Decreto Ejecutivo) detonan transaccionalmente un proceso de actualización sobre los expedientes del módulo Pensionados / Retroactivos.
\end{tcolorbox}

\subsection*{Causas de Retroactivo}

\begin{longtable}{p{4cm}p{8cm}}
\toprule
\textbf{Causa} & \textbf{Descripción} \\
\midrule
\endhead
\bottomrule
\caption{Causas de retroactivo}
\label{tab:causas-retroactivo}
\endlastfoot

Homologación Salarial & Actualización de tabuladores por decreto \\
Nuevos Conceptos & Incorporación de beneficios adicionales \\
Corrección de Cálculo & Rectificación de errores históricos \\
Cambio de Categoría & Reclasificación del pensionado \\
Sentencia Judicial & Orden de pago retroactivo \\
\end{longtable}

\section{Calculadora de Pensiones}

La calculadora permite simular el monto de pensión estimado:

\begin{tcolorbox}[
    enhanced,
    colback=white,
    colframe=sandraDark,
    width=0.85\textwidth,
    halign=center
]
\begin{tabular}{p{5cm}rp{4cm}}
\toprule
\textbf{Parámetro} & \textbf{Valor} & \textbf{Unidad} \\
\midrule
Grado al Retiro & Coronel & \\
Años de Servicio & 25 & Años \\
Fecha de Ingreso & 01/03/2001 & \\
Fecha de Retiro & 28/02/2026 & \\
\midrule
\textbf{Pensión Estimada} & \textbf{\$\$8.500.000,00} & Mensual \\
\bottomrule
\end{tabular}
\end{tcolorbox}

\subsection*{Factores de Cálculo}

La calculadora considera los siguientes factores:

\begin{itemize}[leftmargin=*]
    \item \textbf{Grado}: Nivel jerárquico al momento del retiro
    \item \textbf{Años de Servicio}: Tiempo total de cotización
    \item \textbf{Último Sueldo}: Remuneración base final
    \item \textbf{Componentes}: Bonificaciones aplicables
    \item \textbf{Beneficiarios}: Dependientes registrados
\end{itemize}

\section{Estados del Cálculo}

\begin{longtable}{p{3cm}p{9cm}}
\toprule
\textbf{Estado} & \textbf{Descripción} \\
\midrule
\endhead
\bottomrule
\caption{Estados de cálculo de nómina}
\label{tab:estados-calculo}
\endlastfoot

Preparado & Cálculo generado, en revisión \\
Aprobado & Cálculo validado por supervisor \\
Rechazado & Cálculo con errores, requiere corrección \\
Procesado & Incluido en nómina para pago \\
Pagado & Orden de pago ejecutada \\
\end{longtable}

\section{Consulta de Pagos}

El pensionado puede consultar el estado de sus pagos:

\begin{enumerate}[leftmargin=*]
    \item Seleccionar \textbf{Movimientos} en el menú
    \item Filtrar por período de consulta
    \item Verificar estado de cada orden
    \item Descargar comprobante de pago
\end{enumerate}

\begin{tcolorbox}[
    enhanced,
    colback=sandraSoft,
    colframe=sandraDark,
    width=0.9\textwidth,
    halign=center
]
\small
\textbf{Estados de Pago}:
\begin{itemize}
    \item \textbf{Pendiente}: En cola de procesamiento
    \item \textbf{Ejecutado}: Pago realizado exitosamente
    \item \textbf{Fallido}: Error en la transferencia
    \item \textbf{Reintentando}: Intento de pago en curso
\end{itemize}
\end{tcolorbox}

\section{Reportes del Módulo}

\begin{longtable}{p{4cm}p{8cm}}
\toprule
\textbf{Reporte} & \textbf{Contenido} \\
\midrule
\endhead
\bottomrule
\caption{Reportes del módulo de nómina}
\label{tab:reportes-nomina}
\endlastfoot

Liquidación de Nómina & Detalle completo del cálculo mensual \\
Constancia de Pensión & Certificado oficial del monto pensionario \\
Historial de Pagos & Registro de todos los pagos realizados \\
Proyección Anual & Estimación de ingresos por pensionado \\
Resumen por Grado & Consolidado de pensionados por jerarquía \\
\end{longtable}

\begin{tcolorbox}[
    enhanced,
    colback=sandraSoft,
    colframe=sandraGreen,
    width=0.9\textwidth,
    halign=center
]
\small
\textbf{Importante}: La calculadora proporciona una estimación referencial basada en los datos ingresados. El monto final de la pensión está sujeto a validación por parte del IPSFANB y puede variar según la normativa vigente al momento del retiro.
\end{tcolorbox}

\newpage
